% Comprehensive paper over the series (TN Holdings), first draft, home, 2026-11-28 (owner 12:47 JST "進めて").
% Built from notes/synthesis_2026-11-28/outline.md, papers_read.md, results_ab.md and results_a2_d.md; every number is
% quoted from the paper or record named beside it and is to be checked one by one before release. A try, kept apart
% from paper 12 (notes/paper12/, unchanged).
\documentclass[10pt,a4paper]{article}
\usepackage[margin=2.2cm]{geometry}
\usepackage{amsmath,amssymb}
\usepackage{booktabs}
\usepackage{array}
\usepackage{graphicx}
\usepackage[hidelinks]{hyperref}
\usepackage{float}
\input{glyphtounicode}
\pdfgentounicode=1
\setlength{\parskip}{4pt}
\setlength{\parindent}{0pt}
\sloppy
\newcommand{\pending}[1]{\textbf{[PENDING: #1]}}
\newcommand{\ket}[1]{|#1\rangle}
\newcolumntype{R}[1]{>{\raggedright\arraybackslash}p{#1}}

\title{What a User-Side Layer Can Know Without the Calibration Map:\\ Declared Tests on IBM Heron Devices, 2026}
\author{TN Holdings\\ \small contact: \texttt{love.os.architect@proton.me}}
\date{Draft 1, 28 November 2026}

\begin{document}
\maketitle

\begin{abstract}
A user of a cloud quantum computer sees the device through two numbers it is given: a calibration map that is hours
old when it is read, and results that error mitigation has already changed. We ask what a layer on the user's side can
find out and do with neither: it places circuits from the coupling graph and the failed elements only, applies no
mitigation, puts its checks into the user's own job, and declares every test in git before the data exist. Eleven
preprints (September to November 2026) and the records after them report its pieces; here we test the whole on the
committed records with five analyses declared before they were run. (A) Of 287 predictions declared from 8 October to
27 November, 158 held and 78 failed (15 not read, 36 unclear); about a third of the decided ones failed in every period
until the series declared only predictions an earlier device record grounded, after which none of 17 decided did; the
declared noise models and expected numbers failed most often (20 of 35). (B) Paper 9's line, that a raw result reads
when the share of the exact pattern it keeps is at least $1-1/\sqrt2$, held on 94 of 96 later units, all far above the
line, so the region near it was not tested. (C) Against IBM's own calibration over 191 jobs of the record, the
qubits our in-job mirrors found weak sat on IBM's two-qubit-gate column at 2.29 times chance, while IBM marked no
column for 291 of 721 of them; the map behind a job was more than a day old for 17 of 67 jobs on one device. (D) The
same circuit moved beyond shot noise from day to day in 14 of 16 units, and copies inside one job in none of 9. (E) All
98 re-scoring entries of the record give the committed numbers again. With these, the series' device results read as
follows: raw, unmitigated circuits of up to 153 qubits came within 0.05 of sealed exact answers, and a 91-site
dual-unitary chain kept its decay rates within 0.0037 of theory to 2,205 CZ; where results did not read, the record
says so. No quantum advantage is claimed: every answer here is computed exactly on a classical machine.
\end{abstract}

\section{Introduction}\label{sec:intro}
Two habits stand between a user of today's quantum processors and what the device did. Placement reads the day's
calibration (error rates per qubit and coupler), which ages as soon as it is read and selects in its own favour; and
results are error-mitigated, which changes the number the user sees by a model the user cannot audit. Both help a
benchmark. Neither tells the user how far a given result can be trusted on a given day.

This series took the other position from the start: a layer that sees only the coupling graph, the basis and the
elements the device marks as failed; that puts its own checks, made of the user's circuit, into the same job; that
corrects nothing; and that writes down what it expects before the device runs. Its pieces are reported in eleven
preprints: compiler work that showed where search-based layout fails and where formula-built schedules suffice
(papers 1 to 4, \cite{p1,p2,p3,p4}); placement without the map and the readable depth it gives (papers 5 and 6,
\cite{p5,p6}); an operating layer, psf-os, that carries user-written circuits to the device (papers 7 to 9,
\cite{p7,p8,p9}); and real-time control (papers 10 and 11, \cite{p10,p11}).

This paper does not review them. It asks one question of all their records at once, and answers it with five analyses
(A to E) declared in the record (\texttt{notes/synthesis\_2026-11-28/outline.md}, commit 423b8f3d) before each was run:
how often the series' own predictions held (A); whether the line of paper 9 holds on results that came after it (B);
what IBM's calibration saw of what our checks found (C); how far the same circuit moves from day to day (D); and
whether every recorded number can still be reproduced (E). Sections~\ref{sec:reads} and~\ref{sec:checks} set the
series' device results beside these analyses.

\section{Method common to the series}\label{sec:method}
\textbf{Placement.} A ring by a formula on a lattice cycle (papers 4 to 6), a two-dimensional pattern by translation
of its own edges onto the device's couplers (papers 7 to 9; a pattern that does not fit is refused, not routed), a
chain by a fixed rule round the failed elements (paper 8). No error rate is read. Where this is weaker than search, the
record says so: on a trained classifier formula placement could not be told apart from Qiskit's search
(paper 3)~\cite{p3}, and the three-dimensional prototype needs several times what a search-based router needs
(paper 8)~\cite{p8}.

\textbf{Checks in the same job.} A floor (the user's circuit with its non-diagonal one-qubit gates removed: the same CZ,
ideal $\langle Z\rangle = 1$), a mirror (the compiled circuit and its inverse), a ring sentinel, and, from 23 November,
rulers built from the job's own circuits (string mirrors and readout pairs). They are reported beside a result and never
used to correct it.

\textbf{Verdicts.} A result READS when its mean squared error is at most half the mean square of the exact answer
(papers 7 to 9); AGREES when its difference from a reference is inside a stated width (two standard errors, or the
Student $t$ quantile for $K$ pairs); a real-time result is judged against a correction ceiling (paper 10). Every
answer is computed exactly and sealed by hash before the run.

\textbf{Declarations and records.} Each test's question, circuits, verdict words and predictions are committed before
the device record of its run is fetched; raw counts, job identifiers and the scorers are committed after. From
18 November an independent re-count of each night was made by a separate assistant with separately written code
(Section~\ref{sec:limits} says what kind of independence this is).

\section{What reads (the series' results)}\label{sec:reads}
Table~\ref{tab:reads} gathers the device results with an exact answer, each with the number a benchmark would quote and
the number a user of the device needs.

\begin{table}[H]\centering\footnotesize
\setlength{\tabcolsep}{3pt}
\begin{tabular}{R{2.9cm}R{1.6cm}R{5.2cm}R{5.6cm}}
\toprule
problem & paper, devices & benchmark number & user's number \\
\midrule
22 practice circuits written in plain Qiskit (33 to 153 qubits) & 7; marrakesh, fez & 22 / 22 READ on both; 19 and 18 AGREE; 2D kicked Ising at $\theta = 0.8, 1.0$ within 0.05 to $t = 8$ (worst on fez 0.023) & $\theta = 0.4$ misses by $-0.04$ to $-0.07$ on both; the 33-qubit chain's signal site falls behind with depth (0.30 against 0.41 at $t = 16$); one day \\
dual-unitary kicked-Ising chain, $N = 51, 71, 91$ & 7, 8; marrakesh, fez & decay rates within 0.0045, 0.0017, 0.0037 of $\cos 2h$; every depth to $t = 24$ readable (2,205 CZ at $N = 91$) & $h = 0.15$ sat $+0.0037$ above $\cos 2h$ on both devices \\
3D lattices (kicked Ising $4\times4\times3$; Heisenberg $4\times2\times2$) & 8; marrakesh, fez & Ising read at every angle and depth; Heisenberg read & a light-cone test failed as predicted (every front late); a layout with 46\% fewer CZ came farther from the answer on $4\times3\times2$ \\
IBM's 2023 utility observables; Fermi--Hubbard $U = 4$ & 9; fez & $\langle Z_{106}\rangle$ read raw to 8 steps; Hubbard read at every size from 40 to 120 qubits at 4 Trotter steps & Hubbard read at neither size run at 6 steps \\
kicked Ising ring, conserved charge & 5; marrakesh, kingston & readable depth $L^* = 56$--57 (marrakesh) and 81--90 steps (kingston) against 15--26 for the folded ring and blind search; 71 wins, 11 ties, no loss in 144 comparisons & the declared noise model did not hold; a run of 18 sites on kingston's $n = 40$ ring was dead while the calibration showed nothing \\
\bottomrule
\end{tabular}
\caption{The series' device results with an exact answer, as published~\cite{p5,p7,p8,p9}. Raw: no mitigation, no
post-selection, every shot counted.}\label{tab:reads}
\end{table}

\textbf{B: paper 9's line, after paper 9.} Paper 9 split the error exactly as $\mathrm{MSE}/\overline{e^2} = (1-r)^2 +
\text{rest}$, with $r = \sum ze/\sum e^2$ the share of the exact pattern kept; $r < 0.293$ means ``does not read'' by
construction, and the measured statement was the converse: 160 of the 166 results with $r \ge 0.293$ read, on a table
assembled after the results~\cite{p9}. The results recorded after 16 November that are judged by the same reading are
the Ohm window and ruler nights (23 to 27 November), raw $M$ and raw $W$ per cell: 48 cells, 96 units, each scored by
its own declared scorer (the ratios equal the recorded ones exactly). Declared reading: the line holds if the share that
reads above it is at least $0.964$ minus its binomial 2 SE, $0.935$. It holds: 94 of 96 (0.979). The two that do not
are one cell (A2 $t = 4$, 26 November, ratio 0.529 and 0.526), where, as in paper 9's six, the error that is not a
shrink decided. But every unit sits at $r = 0.94$ to $1.58$; at $t = 3$ and 4 all 60 have $r > 1$ (the device's error
adds to the quantity rather than shrinking it). The test therefore says nothing about the region near 0.293, where
paper 9's misses sat.

\section{What the checks see}\label{sec:checks}
\textbf{Defects the calibration does not show.} On ibm\_kingston a run of 18 sites on the rule's $n = 40$ ring was near
random or inverted in three runs of the byte-identical circuit over three hours, while the calibration showed nothing
there; gate twirling turned the inverted part random, so the cause is on the device~\cite{p5}. A layer that scanned the
rings it would use, listed the dead paths and placed around them restored the ring from 22 to 40 usable sites on three
days (6 of 6 declared comparisons WIDER) and widened rings on two more devices~\cite{p6}.

\textbf{C: against IBM's own calibration.} For the 191 jobs of the record whose submissions name their device and
time, the calibration snapshot that applied when each job ran was fetched and sealed; a scoring memo was committed before the seal was opened, and the calibration
numbers enter no verdict, placement or correction (\texttt{notes/calib\_compare\_2026-11-26/}). Table~\ref{tab:calib}
gives the hit rates. Our mirror findings sat on IBM's two-qubit-gate (CZ) column at 2.29 times chance, against 0.71 to
2.17 on the per-qubit columns; IBM marked no column at all for 291 of 721 of them, and one fez qubit we found 68 times
carried no IBM mark on any. The map behind a job was a median 0.41 to 0.63 hours old, but more than a day old for 17
of 67 jobs on ibm\_marrakesh (largest 93 hours); a qubit's readout error moved by a median factor of 3.30 to 7.81
between its best and worst day.

\begin{table}[H]\centering\footnotesize
\begin{tabular}{lrrrrr}
\toprule
our finding (sets) & readout error & $T_1$ & $T_2$ & CZ & any column \\
\midrule
mirror low, 721 (409) & 0.279 / 0.128 & 0.207 / 0.128 & 0.090 / 0.128 & \textbf{0.488 / 0.214} & 0.596 / 0.441 \\
scan listed, 539 (62) & 0.130 / 0.119 & 0.121 / 0.118 & 0.158 / 0.118 & \textbf{0.336 / 0.194} & 0.583 / 0.458 \\
sentinel dead, 794 (109) & 0.203 / 0.127 & 0.083 / 0.125 & 0.120 / 0.125 & 0.280 / 0.216 & 0.542 / 0.460 \\
\bottomrule
\end{tabular}
\caption{IBM's marks on what our in-job checks found, as hit rate / chance (the share of the set IBM marks on that
column), over 191 jobs on ibm\_fez, ibm\_marrakesh and ibm\_kingston (\texttt{notes/calib\_compare\_2026-11-26/result\_2026-11-27.md}).
IBM misses (no column): 291 of 721, 225 of 539, 364 of 794.}\label{tab:calib}
\end{table}

\textbf{Rulers inside the circuit.} From 23 November the window of an Ohm-type circuit was divided by rulers carried in
the same job: the circuit's own total, mirrors of the circuit and of its strings, and readout pairs (full account in
paper 12, in preparation). Putting the readout's half back into the string mirror's ruler brought the window to the
exact value within the $t$ width in four of six cells at three and four cycles on ibm\_fez (26 November) and five of six
on ibm\_kingston (27 November); the readout alone did not in any of the six on either night. On ibm\_fez every ruler
left a positive residue at three cycles on three nights; on ibm\_kingston the readout-corrected ruler left a residue beyond
its width in one arm of three. A per-qubit record predicted the next night: two named qubits stayed below the declared mirror threshold for a
third night (held), while the next such prediction could not be read when the device rule chose another device.

\textbf{D: the same circuit from day to day.} Table~\ref{tab:spread} compares how far one circuit's raw value moves
between days with its shot error (the counting fixed before the run). Across days the raw value moved beyond shot
noise in 14 of 16 units; copies of a circuit inside one job in none of 9. On the two fez ruler nights the raw window
moved in eight of nine cells, while the window divided by the in-job string mirror moved within its error in all nine
(that error also carries the mirrors' shot noise, so it is the wider of the two).

\begin{table}[H]\centering\footnotesize
\begin{tabular}{R{6.4cm}lrl}
\toprule
unit & days & SD / median shot SE & beyond shot noise ($p < 0.05$) \\
\midrule
fixed daily problem FX-A, ibm\_marrakesh & 5 & 2.71 & yes \\
fixed daily problem FX-A, ibm\_fez & 3 & 1.03 & no \\
Ohm ruler nights, raw $W$ per cell (25, 26 Nov.) & 2 & 1.1--4.4 & 8 of 9 \\
Ohm ruler nights, $W$ divided by the string-mirror ruler & 2 & 0.05--1.7 & 0 of 9 \\
balanced RUS chain per length (17, 18 Nov.; a recalibration between) & 2 & 12--24 & 5 of 5 \\
\midrule
copies inside one job (20 Nov.) & 2 copies & 0.13--1.8 & 0 of 9 \\
\bottomrule
\end{tabular}
\caption{The same circuit, as sent gate for gate, from day to day (\texttt{notes/synthesis\_2026-11-28/results\_a2\_d.md}).}\label{tab:spread}
\end{table}

\section{Control during the run}\label{sec:control}
Papers 10 and 11 asked whether circuits that measure in the middle and act at once get closer to the exact answer than
any correction of the same circuit without the actions can (the correction ceiling). Chains of repeat-until-success
neurons crossed it with a margin growing per neuron, and iterative phase estimation crossed it (0.903 against 0.471 at
$n = 4$)~\cite{p10}. What decided the margin was where qubits wait: a waiting qubit in $\ket{1}$ turned to $\ket{0}$ for
the wait kept $+0.029$ to $+0.052$ more on three devices, and phase estimation built that way stayed beyond its ceiling
to eight bits (0.424 against 0.201)~\cite{p11}. Rule 4 of paper 11, that the same circuits drift between jobs while
copies inside one job agree, is what D finds over days.

\section{The record}\label{sec:record}
\textbf{A: the predictions.} Every prediction with a stated expected outcome declared in the notebook from 8 October to
27 November is one row of a ledger (\texttt{ledger\_a.json}, \texttt{ledger\_a2.json}), with its verdict copied from the
record and its place cited (Table~\ref{tab:ledger}). The count for 16 to 27 November equals the running total the
nights' records keep (75: 54 held, 20 failed, 1 not read). About a third of the decided predictions failed in each
period until 19 November; from 20 November a prediction was declared only where an earlier device record grounded it,
and none of the 17 decided failed. That second number is a change in what was predicted, not evidence of a better model
of the device: the declared noise models, volumes and expected numbers failed in 20 of 35 decided cases. The compiler
work of papers 1 and 2 kept another ledger, which counts correctness checks as hypotheses (1,207 rows, 105 failed); it is
on another unit and is not added here.

\begin{table}[H]\centering\footnotesize
\begin{tabular}{lrrrrrr}
\toprule
period & declared & held & failed & not read & unclear & failed / decided \\
\midrule
8 to 31 October & 139 & 70 & 38 & 9 & 22 & 0.35 \\
1 to 15 November & 73 & 34 & 20 & 5 & 14 & 0.37 \\
16 to 19 November (paper 10's nights) & 57 & 37 & 20 & 0 & 0 & 0.35 \\
20 to 27 November (grounded predictions only) & 18 & 17 & 0 & 1 & 0 & 0.00 \\
\midrule
all & 287 & 158 & 78 & 15 & 36 & 0.33 \\
\bottomrule
\end{tabular}
\caption{Declared predictions and their recorded verdicts (\texttt{notes/synthesis\_2026-11-28/results\_a2\_d.md}).
``Unclear'': the record says ambiguous or gives no verdict for that prediction; each has its reason in the ledger.}\label{tab:ledger}
\end{table}

\textbf{E: re-scoring.} The record's script \texttt{reproduce.py} re-scores each committed device result from its raw
counts with its own scorer; run over all 98 entries on 28 November with psf-os 0.3.3, it returned the committed numbers
in 98 of 98 (to $10^{-9}$). The census covers the scorers; rebuilding the circuits and comparing their gates was not
part of it.

\section{Discussion}\label{sec:discussion}
\textbf{What the layer can say.} For each result it can say whether the result reads and how much of the exact pattern
the device kept, from checks in the same job, without a calibration map (Section~\ref{sec:reads}; B). It can find
defects the calibration does not show and place around them (papers 5 and 6), and its in-job mirrors point at the
same couplers IBM's CZ column marks, at about twice chance, while IBM marks no column at all for about two in five of
the qubits those mirrors find weak (291 of 721; C). Its records can be re-scored by anyone (E), failures included (A).

\textbf{What it cannot do.} It does not make a device better: where the shape of a result is wrong, dividing by a measured
share does not undo it (paper 9), and the rulers of November leave a residue at three cycles on one device that none of
them carries. It does not beat search everywhere: on a trained classifier it could not be told apart from Qiskit's
search, and the three-dimensional prototype costs several times a router (paper 3, paper 8). It does not hold a device
still: the same circuit moved beyond shot noise from day to day in 14 of 16 units (D).

\textbf{What follows for a user.} Compare inside one job: copies there agreed in 9 of 9 units while the same circuit on
different days did not, and a ruler carried in the same job took up most of the day-to-day movement on the one pair of
nights where both were measured (D; the ruler's wider error is part of that). Read a result by the share it keeps and by
its shape, not by a fixed absolute width (paper 7's bias in both directions). Declare predictions only where an earlier
record grounds them, and expect about a third of the rest to fail (A): the ledger gives readers of declared hardware tests
a base rate for how often a stated expectation about a cloud device holds, which we did not find elsewhere.
\pending{check the last clause against the literature before release, or drop it.}

\section{Limits}\label{sec:limits}
\begin{itemize}
\item One device family (IBM heavy-hex; Heron r2 for every device result here), three devices, two months; most single
results are one device and one night.
\item Every answer is classically computable; no claim about what a device can do beyond such problems is made.
\item Tables assembled after their data are named as such: paper 9's line (tested out of sample in B, but far from the
line), the kinds in A (by hand up to 15 November, by fixed rules after).
\item The analyses of this paper were declared before they were run, but by the same team that made the records.
\item The work was done by one person with AI assistants (Claude, by Anthropic): drafting code, declarations and
papers, and the independent re-counts of 18 to 27 November, which were made by a separate assistant instance with
separately written code reading only the raw counts, not by another human team.
\item The calibration comparison is beside the record: it was read after every verdict and changes none.
\end{itemize}

\section*{Data and code}
Records, declarations, raw counts and the analyses of this paper: \url{https://github.com/TN-Holdings-LLC/psf-os-lab-notebook}
(\texttt{notes/synthesis\_2026-11-28/}); the operating layer: \url{https://github.com/TN-Holdings-LLC/psf-os}
(0.3.3, doi:10.5281/zenodo.24496267); the compiler: \url{https://github.com/TN-Holdings-LLC/psf-zero}.

\section*{Acknowledgments}
We acknowledge the use of IBM Quantum services for this work. The views expressed are those of the authors, and do not
reflect the official policy or position of IBM or the IBM Quantum team. The device runs used the IBM Quantum Open Plan
and a Pay-As-You-Go account. \pending{confirm the plan wording against paper 11.}

\begin{thebibliography}{99}\footnotesize
\bibitem{p1} TN Holdings, Ordering Sensitivity in Subgraph-Isomorphism Layout Search, v2 (2026), doi:10.5281/zenodo.22977930.
\bibitem{p2} TN Holdings, PSF-Zero: An Analytic Two-Qubit Gate Synthesizer Combined with a Verified, Ordering-Aware Layout Search, v2 (2026), doi:10.5281/zenodo.22978090.
\bibitem{p3} TN Holdings, When a Quantum Compiler Looks Like a Better Learner: Confounds in Quantum Machine Learning Comparisons, v1.3 (2026), doi:10.5281/zenodo.23826596.
\bibitem{p4} TN Holdings, Folding a Ring onto Heavy-Hex: Formula-Built Schedules with Provable Gate Counts, and What Their Signal Shows on Twins and Devices, v1.3 (2026), doi:10.5281/zenodo.23826764.
\bibitem{p5} TN Holdings, Readable Depth Without the Calibration Map: Formula Placement on Heavy-Hex Devices, and a Defect the Calibration Does Not Show, v1.1 (2026), doi:10.5281/zenodo.23827536.
\bibitem{p6} TN Holdings, What Sets the Readable Depth of a Placed Ring, and How Far a Self-Checking Layer Pushes It, v1.0 (2026), doi:10.5281/zenodo.23883938.
\bibitem{p7} TN Holdings, Exact Answers, No Map, No Mitigation: User-Written Circuits through a Placement-Only OS on IBM Heron r2 Devices, v1 (2026), doi:10.5281/zenodo.24046673.
\bibitem{p8} TN Holdings, Readable Quantum Dynamics Without Error Mitigation: Dual-Unitary Chains up to 91 Qubits and 3D Lattice Tests on IBM Heron Devices, v1.0 (2026), doi:10.5281/zenodo.24136512.
\bibitem{p9} TN Holdings, What Makes Raw Quantum Dynamics Readable? The Share of Signal Kept, Across 179 Results on IBM Heron, v1 (2026), doi:10.5281/zenodo.24278888.
\bibitem{p10} TN Holdings, Decisions No Correction of the Uncontrolled Circuit Can Reach: Real-Time Control on IBM Heron, Declared Before the Data, v1 (2026), doi:10.5281/zenodo.24367001.
\bibitem{p11} TN Holdings, Where a Qubit Waits: Rules for Real-Time Control from Declared Device Tests on IBM Heron, v1 (2026), doi:10.5281/zenodo.24496368.
\end{thebibliography}
\end{document}
